\chapter{引言}

\section{研究背景}

\subsection{应用背景}
过去三十年里，无人化装备从最初的``遥控式''一步步向更自主的方向进化，移动平台上如何高精度地识别和瞄准固定目标，也自然成了工程和研究关注的焦点。军事需求是最早的驱动力：上世纪90年代，美军在``猎人''无人机上测试光电瞄准系统，开创了以视觉为核心的``平台动、目标静''瞄准模式；到了2010年以后，北约和主流军工企业把``视觉伺服加惯性稳定''写进了移动平台火控的设计准则，这个方向才算真正从实验室走向了工程实用\cite{Hutchinson1996visualservo}。

国防之外，民用和商用场景同样迫切需要轻便、便宜又好部署的移动瞄准与跟踪方案。安防监控、交通监测、无人机巡检、精细农业（比如病虫害定点喷洒）、影视拍摄等，哪一样都要求在有限功耗和重量下实现够用的指向精度和鲁棒的跟踪。拿城乡安防来说，街道上的移动巡检车要在各种光照和遮挡条件里持续盯住目标；农业喷洒场景下，设备成本低、续航长是硬约束，根本用不起大功耗传感器阵列，轻量级的视觉加惯导方案自然成了首选。

落到工程实现上，这类系统躲不开一堆约束：传感器和处理器必须在很小的尺寸、重量和功耗预算里协同工作；时间同步和通信延迟直接卡着闭环控制的脖子；复杂环境下还得扛住遮挡、光照变化和多目标干扰。正是这些现实困难，让近几年的研究不再只盯着检测精度，而是更多关注系统级的时延、能耗和可复现性，也给面向教学和工程验证的低成本方案提供了存在的理由。
\subsection{技术背景}
从技术角度看，移动平台瞄准系统要跑起来，离不开视觉传感、惯性测量、实时嵌入式计算和机电稳定（云台）这几个方向的交叉。国内虽然起步晚一些，但2015年以来陆续有几款车载光电转台投入使用，说明在小型化陀螺稳定和实时图像跟踪上已经有了工程级的突破。平时衡量一套轻量化方案行不行，常用``千元成本、5\,W功耗、500\,g重量''这类指标——也恰恰是这些约束在推动学术研究往工程实现上靠。

几个关键技术维度值得展开说。第一是传感器怎么选：工业相机、激光雷达、民用USB摄像头、低成本六轴IMU……选什么直接决定了精度、体积、功耗和成本之间的取舍。第二是实时处理和算法：传统方法比如HOG+SVM在资源受限场景下性价比不错，而轻量化神经网络（MobileNet、Tiny-YOLO这类）最近几年靠模型压缩和量化，也慢慢能在嵌入式平台上跑出可用帧率了。第三是时间同步：高频IMU和低频视觉要做紧耦合，精确时间戳和低延迟通信是前提，否则前馈补偿和闭环控制的效果会大打折扣。第四是机械和控制：云台的机械结构好不好、减震和电机驱动策略到不到位，最后都反映在指向精度上。

软件架构上，嵌入式平台（比如STM32H7）配合上位机（Linux视觉端）的分工方式已经比较普遍了：上位机干视觉推理和图像处理的重活，主控板管IMU采样、闭环控制和实时执行，这样既保证了性能，又不至于把主控端算力撑爆。常见的工程做法还包括用DMA和中断来保证IMU采样的高精度，走RS485这类差分总线跟云台电机可靠通信，以及在控制环路里加前馈和预测补偿来弥补视觉反馈滞后的亏。

总的来看，技术背景有两条并行的线索：一条是追求极致精度，还是靠高端传感器和大算力平台；另一条是为了满足成本和功耗约束，把研究重心往算法和系统协同优化上挪——比如轻量模型、传感器融合、时序对齐这些方向。后一条路恰恰给本科和工程验证级别的研究留出了可行的空间。
\section{研究目的及意义}

\subsection{研究目的}
本课题设计并实现基于移动平台的固定目标识别及瞄准系统，采用``纯视觉+IMU''紧耦合的轻量化方案，在满足成本约束的前提下，实现以下具体目标：
\begin{enumerate}
\item 构建硬件平台与通信基础：搭建以STM32H7为核心的小型化云台系统，集成USB摄像头与六轴IMU，实现云台的实时远程控制；
\item 实现环境感知与角速度采集：完成高频（≥100 Hz）运动基座状态信息采集与处理；
\item 完成视觉跟踪与像素闭环：实现25 FPS连续跟踪，瞄准线稳态误差控制在0.8 mrad以内；
\item 集成IMU前馈与振动抑制：在像素闭环基础上引入IMU补偿，对10\textasciitilde{}200 Hz频段进行实时补偿。
\end{enumerate}

\subsection{研究意义}
这项工作在学术层面填补了国内在轻量化动基座瞄准方案上的一些空白，在应用层面则为无人系统提供了一条可复现、低成本的参考路径。搭起``纯视觉加低成本''这个设计框架之后，未来如果想接入激光雷达或者事件相机，硬件和算法的接口也都预留好了。
\section{研究现状及设计思路}

\subsection{国外研究现状}
国外在这方面起步早、投入大，基本上走的是``高精度、高成本、重平台''的路子。传感器方案上，激光雷达加视觉再加惯导的三元融合几乎成了标配。

举个例子，卡内基梅隆大学2022年公开的实验里，他们的无人地面车辆装了128线激光雷达、工业级可见光相机（1920×1440@30\,Hz）和光纤陀螺，组成了一套三位一体的传感器系统。靠Transformer检测头\cite{Ren2015FasterRCNN}做到30\,Hz实时识别，能在复杂城区环境里精确定位固定目标，像素误差控制在2个像素以内，折合角度精度大约0.3\,mrad。但问题也很明显：整套系统成本超过2万美元，功耗45\,W，对移动平台的续航是个很大的拖累。这种方案放到本科教学层面，想做到``低成本、易复现''基本不现实。

德国的弗劳恩霍夫应用研究所（Fraunhofer IAO）走得更远，把FPGA和GPU的异构计算架构引进了瞄准闭环。GPU上跑Faster R-CNN做目标检测，FPGA上并行执行IMU前馈补偿，利用异构计算的高吞吐把定位误差压到了1\,cm量级（对应0.05\,mrad的角误差）。但代价也不小：系统重6\,kg，供电得24\,V、10\,A，单是陀螺稳定分系统的成本就超过5000元，想在5\,kg以下的微型无人车上部署，几乎不可能。

传感器层面，国外普遍用工业级高精度IMU（漂移率$<1^\circ$/h）和光纤陀螺（精度$<0.01^\circ$/s），比民用级传感器贵出一个数量级还不止。Xsens的MTi-630 AHRS要8000元以上，光纤陀螺类产品更是数万元起步。靠堆传感器换取高精度，代价实在太大了。

总的来看，国外研究强调``多传感器融合加大算力换高精度''，靠物理层的传感器冗余和算法层的计算冗余把性能指标推到极致。精度上确实遥遥领先\cite{Redmon2016YOLO}，但这种技术路线对资源受限的场景——比如微型无人车、教学实验平台——天然就水土不服。
\subsection{国内研究现状}
受成本约束影响，国内研究走的路子要务实得多。大多数团队选择的是``纯视觉加惯导''的轻量化路线\cite{Jia2000MachineVision}，核心理念是：硬件配置不够，算法设计来凑——在有限的计算资源和传感器精度下，仍然要让瞄准精度达到工程可用的水平。

早期探索方面，南京航空航天大学2013年的硕士论文走在了前面。他们在320×240分辨率的低性能嵌入式平台上，用HOG加SVM分类器\cite{Dalal2005HOG}做出了车载固定目标识别，单发命中精度达到0.8\,mil。这说明传统机器学习在资源受限场景下是可行的。不过，这项工作没能有效解决``边走边打''（也就是动基座打静止目标）过程中平台振动带来的零偏问题——零偏往往在100\,ms量级，足够把瞄准精度毁掉，这个问题后来成了后续研究的一个重点。

到了深度学习时代，齐鲁工业大学2021年的学位论文往前迈了一大步。他们把Faster R-CNN\cite{Ren2015FasterRCNN}的骨干网络换成轻量化的MobileNet，在STM32加OpenMV的架构上跑出了20\,FPS的实时检测，mAP达到了72\%，比HOG+SVM明显好。但一个突出的工程短板也暴露出来：遮挡场景或者强光突变下，目标丢失率仍然有6\%~8\%，系统的鲁棒性还是不够。

电子科技大学2022年的一项嵌入式视觉研究，把这个瓶颈进一步量化了。他们仔细分析了性能数据：就算用轻量化的MobileNet CNN，在640×480的分辨率下，推理过程仍然会吃掉嵌入式GPU 80\%的算力。这种高负载不光导致瞄准延迟超过100\,ms（实时性要求可是10\,ms量级），还会带来明显的功耗上升和热量积累，威胁嵌入式平台的稳定性。

算法层面，国内在IMU数据处理上也有进展。好几个团队研究了IMU零偏校准和温漂补偿，但多数还停在离线标定阶段，实时的自适应标定和在线零偏补偿很少见。至于``像素误差加角速度前馈''这种紧耦合双闭环控制，在国内还处于萌芽期，缺乏系统的理论分析和工程验证。

总的来说，国内轻量化方案在振动抑制、实时性和鲁棒性上，跟国外大概有1到2年的工程化差距。原因主要有三个：一是国内团队在算法创新上积累不够，习惯跟着国外最新论文走，工程适配性没跟上；二是轻量化IMU和高性能嵌入式芯片在国内不太好拿；三是``纯视觉加IMU''紧耦合的完整系统设计和验证，缺少学科交叉的系统研究。
\subsection{存在的主要问题}
综合国内外研究现状，当前存在的主要问题包括：
\begin{enumerate}
\item 国外方案在成本、体积和功耗方面存在极端不平衡，不适用于轻量化应用；
\item 国内轻量化方案在振动抑制、实时性与鲁棒性上与国外存在1\textasciitilde{}2年的工程化差距；
\item ``像素-角速度''一体化设计尚处萌芽阶段，``纯视觉+IMU''紧耦合的轻量化方案缺乏系统研究；
\item 移动平台动态运动对瞄准精度的影响补偿不足，需要建立双闭环控制架构。
\end{enumerate}
\subsection{本文的设计思路}
综合以上分析，本工作确定了以下几条设计原则：
\begin{enumerate}
\item \textbf{架构选择}：走``纯视觉+IMU''紧耦合的路线，不碰激光雷达这种昂贵的传感器，在``千元级成本、5\,W功耗''的框框里做出可接受的精度；
\item \textbf{特征算法}：用HOG、ORB这类传统特征和模板匹配，不依赖GPU，在嵌入式平台上就能跑，保证帧率超过25\,FPS；
\item \textbf{控制架构}：搭一套``像素误差外环加角速度前馈内环''的双闭环，靠IMU数据提前预估图像偏移，把平台机动带来的误差降下来；
\item \textbf{时间同步}：硬件层面让图像和IMU数据的时间戳对齐，给紧耦合打好底。
\end{enumerate}
这样设计，既保证了工程上的可复现性，也为将来技术升级留了余地。
\section{发展趋势}
国内移动平台瞄准系统正在从``多传感器融合''向``视觉-IMU''紧耦合的轻量化方向转变。硬件上，国产微控制器的主频已经突破了480\,MHz，内置FPU和DSP指令集，给``视觉跟踪加IMU前馈''同频运行提供了算力保障。算法上，HOG、ORB这类传统特征不要GPU就能跑出高帧率，而且跟IMU采样频率接近，时间戳对齐也方便。控制上，``像素误差外环、角速度前馈内环''的双闭环已经比较成熟，能把平台机动引起的瞄准线误差压缩掉30\%以上。

可以说，``纯视觉加IMU''紧耦合已经是低成本动基座瞄准的一个明确趋势，也给本科毕业设计提供了一个能落地、能复现的技术窗口。
\section{主要内容和论文组织}
本文围绕移动平台固定目标识别与瞄准问题，设计并实现了一套低成本系统。全文6章，内容安排如下：

\textbf{第1章（引言）}交代课题的研究背景和应用意义，梳理国内外研究现状、存在的问题、设计思路和发展趋势，说明为什么做这个课题、能不能做。

\textbf{第2章（系统总体设计）}给出系统整体方案，明确``Linux视觉端+云台主控+底盘主控''的三节点分布式架构，说明各节点之间的通信协议和数据流、系统工作原理和框图，同时列出功能需求和性能指标约束。

\textbf{第3章（系统硬件设计）}逐一讲解硬件各模块的详细设计，包括STM32H7主控核心电路、HiPNUC六轴IMU采集电路、USB摄像头接入方式、MS4010双轴云台驱动电路和底盘传感器接口，最后给出整机连接关系图。

\textbf{第4章（系统软件设计）}详细描述软件的分层架构和实现细节——IMU角速度采集与处理、视觉目标检测与跟踪、像素误差外环和IMU前馈内环双闭环控制算法、底盘差速驱动与八路灰度循迹控制，以及节点间的通信协议实现。

\textbf{第5章（系统调试与结果分析）}分模块和整机两个层次进行调试，包括图像识别与跟踪效果测试、IMU采样稳定性与补偿性能验证、云台闭环稳态误差测试和底盘循迹精度分析，给出实测结果和相应分析。

\textbf{第6章（总结与展望）}归纳主要工作和创新点，分析系统当前存在的不足，并讨论后续可以深化的方向和技术升级路径。


